\abstract{
We propose PY-IDR, a hierarchical Bayesian nonparametric extension of the Irreproducible Discovery Rate (IDR) framework of \citet{LiBrownHuangBickel2011}. PY-IDR is designed for an arbitrary number of replicates $K \geq 2$, replicate-specific and non-Gaussian marginal score distributions, and an unknown number of reproducible-component dependence regimes. The reproducible copula is modeled as a Pitman--Yor mixture of Gaussian copulas with structured correlation matrices, augmented with $t_5$, Clayton, and Gumbel atoms for planned tail-dependence robustness; replicate marginals are modeled by Bernstein--Dirichlet priors with a hierarchical degree prior. Inference is developed through a Pitman--Yor P\'olya-urn MCMC scheme with auxiliary atoms, NUTS or Metropolis updates on valid correlation parameterizations, and a finite-truncation structured variational alternative for large genomic studies. The theory establishes finite-sieve identifiability under positive lower-orthant dependence and explicit component-separation assumptions, gives a conditional Hellinger contraction statement under stated KL-support and sieve-tail conditions, connects posterior local-idr consistency to Bayes-mFDR control under threshold regularity assumptions, and separates Harris ergodicity from a conjectural geometric-rate strengthening. We also provide an empirical protocol spanning the S1--S4 IDR benchmarks, a heavy-tail scenario S5, and five planned real-data applications: ENCODE 4 CTCF ChIP-seq, ENCODE 4 ATAC-seq, DREAM5, Tabula Muris pseudobulk, and Pan-UK Biobank GWAS. The present manuscript develops the model, inference scheme, theory, implementation roadmap, and evaluation protocol; its figures and result tables are illustrative or placeholder material, and final observed numerical values will be reported only after the companion experimental implementation is run.
}
